% Chapter 14 of Verifiable Intelligence: proofs of training, committed logs, and refereed training.

\protocolchapter{Proofs of Training}{Arguments about how a model was made}{logged-checkpoint spot checks}{C,S,Q,Q,Q,D}{ch:training}
\begin{refsection}[chapters/14-training.bib]

\begin{chaptersummary}
A developer states how a model was made, and an auditor who cannot repeat the run wants to check the statement. The developer commits to its log of checkpoints with a Merkle tree, and the auditor replays a few segments chosen at random. We prove that the commitment binds a developer who cannot find hash collisions, so that a log with a fraction $\varepsilon$ of invalid segments is caught with probability close to $1-(1-\varepsilon)^q$, as in Kilian's argument; that a log can be valid within the tolerance that nondeterministic hardware forces although no training was done; and that two trainers, one of them honest, let a referee settle the exact result by replaying one step. Honest training passes but is enforced only when the specification fixes data, seed, and hyperparameters; the auditor replays a fraction $q/N$ of the run; skipped work is caught, but a wrong model needs only one invalid segment. Left out are the provenance of the data, the deployed model, and the cost of replay at frontier scale.
\end{chaptersummary}

\section{The oversight problem: a training run nobody can repeat}\label{sec:training:problem}

A model's weights determine what it computes, not how it was made, and some questions an overseer must answer are about the making.

\begin{oversight}[a run below the threshold]
A developer publishes the weights of a language model with a statement: training started from the initialization generated by a published seed, used a data set whose fingerprint was published before training began, and ran a published recipe for $200{,}000$ optimizer steps, about $8\times10^{24}$ floating-point operations, below a threshold of $10^{25}$ above which the law attaches obligations. An auditor may read the developer's records and owns a cluster on which it can replay about one percent of the run. It cannot repeat the whole run, and if it could, the replay would not reproduce the weights bit for bit.
\end{oversight}

The scenario is invented; its threshold is the one at which the European Union's AI Act presumes that a general-purpose model has high-impact capabilities \parencite[Art.~51(2)]{EU2024AIAct}. The weights carry no record of their history; records handed over after training could have been written after training; and a replay costs what the training cost and, on other hardware, ends elsewhere, because floating-point additions are not associative and their order on a graphics processor varies from run to run \parencite{JiaEtAl2021,SrivastavaAroraBoneh2024}.

\Textcite{JiaEtAl2021} proposed \term{proof-of-learning}\index{proof-of-learning}: the developer logs a checkpoint every few steps, with the batch indices and hyperparameters, and the verifier replays the segments with the largest updates in each epoch and accepts if each replay ends within a tolerance of the next checkpoint. They argued its security from the growth of entropy during training, and stated as a conjecture that inverting a training run, the natural way to forge a log for a given model, costs at least as much as training;\sidenote{They also test the initialization, with a Kolmogorov--Smirnov test against the claimed distribution, and accept a warm start only with a proof of learning for its initial weights \parencite{JiaEtAl2021}.} attacks followed \parencite{ZhangEtAl2022,FangEtAl2023}. \Textcite{Shavit2023} proposed that chips save snapshots of the weights at random times, and \textcite{ChoiShavitDuvenaud2023} tests of a model's training data, which they call proof-of-training-data. Zero-knowledge proofs of training exist for small models at a large cost to the prover \parencite{GargEtAl2023,AbbaszadehEtAl2024}, and \textcite{SrivastavaAroraBoneh2024} made training replayable bit for bit, so that an auditor who repeats a run can dispute it step by step.

The tools are older. \Textcite{Merkle1987} committed to a long list with a tree of hashes; \textcite{Kilian1992} committed in this way to a probabilistically checkable proof and opened only the bits the verifier asked for, which made the spot checks of \cref{ch:spotchecks} usable against a prover who answers questions;\sidenote{Kilian's paper is about zero-knowledge arguments; the zero-knowledge part is left to \cref{ch:zeroknowledge}.} \textcite{Micali2000} removed the interaction; and \textcite{CanettiRivaRothblum2011} let a client who hands a computation to several servers, one of them honest, find the correct result by bisection. What can an auditor who replays a small fraction of a run learn about how a model was made, and what must the specification fix for the answer to mean anything?

\section{The protocol: commit to the log, replay random segments}\label{sec:training:protocol}

We model a training run as a deterministic computation that the specification describes completely, and the developer's records as a list of checkpoints. All norms are Euclidean.

\begin{definition}[Training specifications and logs]\label[definition]{def:training:spec}\index{training log}
A \term{training specification} fixes a dimension $d$, an initial model $w_0\in\R^d$, \emph{segment maps} $F_1,\dots,F_N\colon\R^d\to\R^d$ and a tolerance $\tau\ge0$; the map $F_j$ runs the $\Delta$ optimizer steps of segment $j$ on the batches the specification assigns to them, with its hyperparameters, so the run has $T=N\Delta$ steps. A \emph{log} from $w_0$ to $\hat w$ is a sequence $L=(W_0,\dots,W_N)$ of checkpoints with $W_0=w_0$ and $W_N=\hat w$; its segment $j$ is \emph{valid} if $\norm{F_j(W_{j-1})-W_j}\le\tau$, and $e(L)$ is the number of its invalid segments. The \emph{honest log} has $W_j=F_j(W_{j-1})$ for every $j$, and a model $\hat w$ is \emph{reachable} if some log from $w_0$ to $\hat w$ has no invalid segment.
\end{definition}

When $\tau=0$, the only reachable model is $F_N\circ\dots\circ F_1(w_0)$, the result of the honest run. Every model has a log with at most one invalid segment, which follows the honest run for $N-1$ segments and then jumps to $\hat w$; so spot checks find skipped work, which leaves many invalid segments, far more easily than a wrong model. Below, $\varepsilon$ is a fraction of invalid segments, not a learning accuracy.\sidenote{As in property testing, the protocol separates a fully valid log from logs with many invalid segments, and promises nothing in between.}

The developer commits to its log with a \term{Merkle tree}\index{Merkle tree} (\cref{def:toolkit:crhf}). A hash function $H$ with outputs of $\lambda$ bits, drawn by its key from a collision-resistant family, is applied to each checkpoint and then to the concatenation of each pair of sibling nodes, up to a root $c$, the \emph{commitment}; the $N+1$ positions are padded with copies of a fixed string to $2^h$ leaves, $h=\lceil\log_2(N+1)\rceil$. To \emph{open} position $i$ is to send $W_i$ and the $h$ siblings $s_1,\dots,s_h$ of the nodes on its path; the receiver computes $x_0=H(W_i)$ and $x_\ell=H(x_{\ell-1}\,\|\,s_\ell)$ or $H(s_\ell\,\|\,x_{\ell-1})$ as the $\ell$th lowest bit of $i$ is $0$ or $1$, and accepts the opening if $x_h=c$. We write $\gamma(t)$ for the largest probability, over the key and its coins, with which an algorithm running in time $t$ finds two different strings with the same hash; collision resistance says that $\gamma(t)$ is negligible in $\lambda$ when $t$ is polynomial in $\lambda$.

\begin{protocol}{Logged-checkpoint spot checks}\label{prot:training:spot}
\item[\msg{V}{P}] A key for the hash family, drawn at random, before training.
\item[\msg{P}{V}] After training, the commitment $c$: the root of the Merkle tree over its log $(W_0,\dots,W_N)$.
\item[\act{V}] Draw $j_1,\dots,j_q$ independently and uniformly from $\{1,\dots,N\}$.
\item[\msg{V}{P}] The indices $j_1,\dots,j_q$.
\item[\msg{P}{V}] Openings of the positions $0$ and $N$ and of $j_i-1$ and $j_i$ for $i=1,\dots,q$.
\item[\act{V}] Check every opening against $c$, check $W_0=w_0$ and $W_N=\hat w$, and for each $i$ replay segment $j_i$ and check $\norm{F_{j_i}(W_{j_i-1})-W_{j_i}}\le\tau$. \acc{} if every check passes; otherwise \rej.
\end{protocol}

The developer $\Prov$ has trained and kept its checkpoints; the auditor $\V$ knows the specification and the published model $\hat w$, and can compute each $F_j$. The key is drawn before the commitment, so that the developer cannot pick a hash function in which it knows collisions,\sidenote{Deployed systems use one fixed function, such as SHA-256, for which collisions exist, since it compresses, but are believed infeasible to find; a keyed family is how the theory states that belief (\cref{def:toolkit:crhf}).} and the checks after it, so that the developer cannot choose what to open after seeing them. The auditor compares a replay with the opened checkpoint, not with its hash, since the two agree only within $\tau$.\sidenote{For this reason \textcite{JiaEtAl2021} did not hash their logs; they mention fuzzy extractors and locality-sensitive hashing as ways to compare hashes approximately, and the difficulty of using them.}

\begin{table}[t]
\centering\small
\begin{tabular}{@{}c c c c c c c c c@{}}
\toprule
segment $j$ & $1$ & $2$ & $3$ & $4$ & $5$ & $6$ & $7$ & $8$\\
\midrule
batch $b_j$ & $8$ & $6$ & $7$ & $2$ & $10$ & $1$ & $6$ & $9$\\
honest log $W_j$ & $4$ & $5$ & $6$ & $4$ & $7$ & $4$ & $5$ & $7$\\
lazy log $W_j$ & $4$ & $5$ & $6$ & $4$ & $4$ & $4$ & $4$ & $4$\\
replay $F_j$ of the lazy $W_{j-1}$ & $4$ & $5$ & $6$ & $4$ & $7$ & $2.5$ & $5$ & $6.5$\\
valid & yes & yes & yes & yes & no & no & no & no\\
\bottomrule
\end{tabular}
\caption{Eight one-step segments from $w_0=0$ with $F_j(w)=(w+b_j)/2$, a step of size $\frac12$ on the loss $\frac12(w-b_j)^2$, and $\tau=0$. A lazy developer stops after four segments and publishes $\hat w=4$; four of its segments are invalid.}\label{tab:training:run}
\end{table}

\paragraph{A worked run.} \Cref{tab:training:run} shows a run with $d=1$, $N=8$ one-step segments, and $\tau=0$: the honest log ends at $7$, and a lazy developer stops after four segments and publishes $\hat w=4$. The nine positions are padded to $16$ leaves, so each opening carries $h=4$ hashes. If the auditor draws $j=(3,7,2)$, the developer opens positions $0,1,2,3,6,7,8$; segment $3$ replays to $(5+7)/2=6=W_3$, but segment $7$ replays to $(4+6)/2=5\ne4$, and the auditor rejects. Three checks miss the four invalid segments with probability $(4/8)^3=1/8$. A developer could defend $\hat w=4$ at less risk, though not with less work, by training for seven segments and jumping from $5$ to $4$ in the last: that log has one invalid segment, and three checks miss the jump with probability $(7/8)^3\approx0.67$.

\begin{dictionary}
the developer $\Prov$ & the developer's training cluster, which keeps checkpoints\\
the auditor $\V$ & a regulator or auditor with trusted hardware for replay\\
the specification $(w_0,F_1,\dots,F_N,\tau)$ & the published seed, data fingerprint, recipe, and length\\
the commitment $c$ & a fingerprint of the whole log, published before the audit\\
opening position $i$ & handing over one checkpoint and the hashes that tie it to the fingerprint\\
an invalid segment & part of the run that was not computed as the record claims\\
the tolerance $\tau$ & the allowance for hardware that does not replay exactly\\
$\hat w$ is reachable & some account of the model by the recipe skips no part of the run
\end{dictionary}

What one had to notice is that a training log is an ordinary proof in the sense of \cref{ch:spotchecks}, a long sequence of steps each checkable by itself. Spot checks of such a proof are weak against a single bad step (\cref{prop:spotchecks:naive}), and a developer, unlike a printed proof, answers questions after they are asked; the commitment removes the second difficulty, as it did for \textcite{Kilian1992}, and leaves the first. And the checks certify a property of the log, while the overseer cares about a property of the world, linked to it only by what the specification fixes.

\section{The theory: binding and spot checks}\label{sec:training:binding}

We show that the root binds the developer to one string, that random checks of a fixed log find invalid segments, and that the two combine against developers who cannot find collisions.

\begin{theorem}[Merkle commitments are binding]\label{thm:training:merkle}\index{commitment!binding}
Let $H$ be drawn from a collision-resistant family, and let $c$ be the root of a Merkle tree with $2^h$ leaves.
\begin{enumerate}
\item From two openings of the same position $i$ to different strings $v\ne v'$, both accepted against $c$, one obtains a collision of $H$ among the at most $2(h+1)$ evaluations that checking them takes. Hence an algorithm that runs in time $t$ outputs a root and two such openings with probability at most $\gamma(t+t_0)$, where $t_0$ is the time to check two openings.
\item Committing takes $2^{h+1}-1$ evaluations of $H$; an opening consists of the string and $h$ hashes of $\lambda$ bits, and checking it takes $h+1$ evaluations.
\end{enumerate}
\end{theorem}

\begin{proof}
The plan is to walk up the two paths together until they meet. Let $x_0,\dots,x_h$ and $x'_0,\dots,x'_h$ be the values computed in checking the two openings, so that $x_h=x'_h=c$, and let $z_\ell$ and $z'_\ell$ be the inputs of the evaluations that produced $x_\ell$ and $x'_\ell$: $z_0=v$, $z'_0=v'$, and for $\ell\ge1$ the concatenation of $x_{\ell-1}$, or $x'_{\ell-1}$, with a sibling, in the order that the $\ell$th bit of $i$ prescribes, the same for both openings. Let $\ell^*$ be the least $\ell$ with $x_\ell=x'_\ell$; it exists because $x_h=x'_h$. The inputs $z_{\ell^*}$ and $z'_{\ell^*}$ differ: for $\ell^*=0$ because $v\ne v'$, and for $\ell^*\ge1$ because $x_{\ell^*-1}\ne x'_{\ell^*-1}$ by the choice of $\ell^*$, and the two occupy the same half of the concatenations. So $(z_{\ell^*},z'_{\ell^*})$ is a collision. An algorithm that outputs the openings and then checks them finds a collision in time $t+t_0$ whenever the openings are accepted and $v\ne v'$, which gives (i). For (ii), the tree has $2^h$ leaves and $2^h-1$ internal nodes, one evaluation each, and a check computes $x_0,\dots,x_h$.
\end{proof}

\begin{proposition}[Spot checks of a committed log]\label[proposition]{prop:training:spotcheck}\index{spot check}
Let a log with $e$ invalid segments among $N$ be fixed before the checks are drawn, and let $\varepsilon=e/N$.
\begin{enumerate}
\item If $q$ checks are drawn independently and uniformly from the $N$ segments, then all of them miss the invalid segments with probability exactly $(1-\varepsilon)^q$, and an invalid segment is found with probability $1-(1-\varepsilon)^q$.
\item If $q\le N$ distinct segments are drawn uniformly, then all miss with probability $\binom{N-e}{q}\big/\binom Nq\le(1-\varepsilon)^q$.
\item If $0<\varepsilon<1$ and $0<\beta<1$, then detection with probability at least $1-\beta$ in case (i) needs $q\ge\ln(1/\beta)/\ln(1/(1-\varepsilon))$, and $q=\lceil\ln(1/\beta)/\varepsilon\rceil$ suffices.
\end{enumerate}
\end{proposition}

\begin{proof}
(i) Each check misses with probability $(N-e)/N=1-\varepsilon$, independently of the others. (ii) is \cref{prop:spotchecks:naive}(ii). (iii) The bound $(1-\varepsilon)^q\le\beta$ is equivalent to the first inequality, and $(1-\varepsilon)^q\le e^{-\varepsilon q}\le\beta$ when $q\ge\ln(1/\beta)/\varepsilon$.
\end{proof}

Against a single invalid segment, $q$ distinct uniform checks detect with probability exactly $q/N$, and no rule of $q$ checks does better against a developer who knows the rule (\cref{prop:spotchecks:naive}(i)). The proposition needs the log to be fixed before the draw. A developer who sends only a root has fixed nothing in the information-theoretic sense, but it has fixed a string in the sense that matters against bounded developers. To say this for logs and for probabilistically checkable proofs at once, call a verifier a \emph{spot-check verifier} if, on input $x$ and with coins $\rho$, it computes at most $\kappa$ positions of a string $\pi$ of $M$ symbols from $x$ and $\rho$ alone, reads the symbols there, and accepts or rejects; its soundness error on $x$ is $\delta(x)=\max_\pi\Prob_\rho[\V\text{ accepts }\pi]$. In its \emph{committed version} the verifier draws a key, the prover sends the root of a Merkle tree over its string, the verifier draws $\rho$, and the prover opens the positions that $\rho$ selects; the verifier rejects if an opening fails its check. Let $t_{\V}$ be the verifier's running time, its checks of the openings included.

\begin{theorem}[Committed spot checks]\label{thm:training:committed}\index{argument}\index{soundness!computational}
Let $\V$ be a spot-check verifier that reads at most $\kappa$ of $M$ positions, and run its committed version with a tree of depth $h=\lceil\log_2M\rceil$.
\begin{enumerate}
\item \emph{Completeness.} If $\V$ accepts a string $\pi$ for every $\rho$, then the prover that commits to $\pi$ and opens it is accepted with probability $1$.
\item \emph{Soundness.} Let $\Prov^*$ run in time $t^*$, and for each key and choice of its coins let $\tilde\pi$ hold at each position the symbol that $\Prov^*$ opens there most often in accepting runs. Over the key, the coins of $\Prov^*$, and $\rho$, the verifier accepts after reading a symbol off $\tilde\pi$ with probability at most $\sqrt{\kappa M\,\gamma(2t^*+2t_{\V})}$, and accepts $x$ with probability at most
\begin{equation}\label{eq:training:committed}
\delta(x)+\sqrt{\kappa M\,\gamma(2t^*+2t_{\V})}.
\end{equation}
\item \emph{Costs.} The prover evaluates $H$ $2^{h+1}-1$ times to commit; the verifier receives the root and at most $\kappa$ openings and evaluates $H$ at most $\kappa(h+1)$ times besides its own computation.
\end{enumerate}
\end{theorem}

\begin{proof}
Completeness holds because honest openings are accepted, and (iii) is \cref{thm:training:merkle}(ii). For soundness, the plan is to fix the prover's coins, define the string it opens most often at each position, bound the runs that agree with that string by $\delta(x)$, and turn the runs that disagree into collisions by running the prover twice on the same coins.

Fix the key and the coins of $\Prov^*$, together $\omega$; then the root and the answers to every $\rho$ are fixed. For a position $p$ and a symbol $v$, let $\alpha_p(v)$ be the probability over $\rho$ that $\V$ reads position $p$, receives an accepted opening of it to $v$, and accepts. Let $\tilde\pi_p$ be a symbol that maximizes $\alpha_p$, let $S_p=\sum_v\alpha_p(v)$, and let $m_p=S_p-\alpha_p(\tilde\pi_p)$, the probability that $\V$ accepts after reading at $p$ a symbol other than $\tilde\pi_p$. An accepting run either reads the symbols of $\tilde\pi$ at all its positions, and then $\V$ accepts $\tilde\pi$ with the same coins, because its positions and decision depend only on $x$, $\rho$, and the symbols read; or it reads some position $p$ off $\tilde\pi$, which by the union bound has probability at most $\sum_pm_p$. Hence
\begin{equation}\label{eq:training:split}
\Prob_\rho[\V\text{ accepts}\mid\omega]\le\Prob_\rho[\V\text{ accepts }\tilde\pi]+\sum_pm_p\le\delta(x)+\sum_pm_p.
\end{equation}

To bound $\sum_pm_p$, let an algorithm $B$, given a key, draw coins for $\Prov^*$, run $\Prov^*$ with them against two independent draws $\rho$ and $\rho'$ of the verifier's coins, and, if the two runs open some position to different symbols with accepted openings, output the collision of \cref{thm:training:merkle}(i); it runs in time $2t^*+2t_{\V}$. For each $p$, the two runs accept and read position $p$ as different symbols with probability
\[\sum_{v\ne v'}\alpha_p(v)\alpha_p(v')=S_p^2-\sum_v\alpha_p(v)^2\ge S_p^2-\alpha_p(\tilde\pi_p)S_p=S_pm_p,\]
since $\sum_v\alpha_p(v)^2\le\alpha_p(\tilde\pi_p)\sum_v\alpha_p(v)$. Two runs that disagree do so at no more than $\kappa$ positions, those the first run reads, so the probability $\gamma_\omega$ that $B$ succeeds with the coins $\omega$ satisfies $\kappa\gamma_\omega\ge\sum_pS_pm_p$. By the Cauchy--Schwarz inequality over the positions with $S_p>0$, and $m_p\le S_p$,
\[\sum_pm_p\le\Bigl(\sum_pS_pm_p\Bigr)^{1/2}\Bigl(\sum_p\frac{m_p}{S_p}\Bigr)^{1/2}\le\sqrt{\kappa\gamma_\omega M}.\]
The mean of a square root is at most the square root of the mean, and $\E_\omega\gamma_\omega$ is the probability that $B$ finds a collision, at most $\gamma(2t^*+2t_{\V})$; so the mean of $\sum_pm_p$ over $\omega$, which bounds the probability of accepting after reading off $\tilde\pi$, is at most $\sqrt{\kappa M\gamma(2t^*+2t_{\V})}$, and averaging \eqref{eq:training:split} over $\omega$ gives \eqref{eq:training:committed}.
\end{proof}

The square root and the factor $M$ are the price of running the prover twice on the same coins, which is called \emph{rewinding}.\sidenote{Rewinding is also how one proves knowledge soundness, that an accepted prover ``knows'' a proof that can be extracted from it; \textcite{Thaler2022} discusses it for Kilian's argument. In the random-oracle model an extractor can read the tree off the prover's queries without rewinding, and the loss becomes additive \parencite{BenSassonChiesaSpooner2016,Thaler2022}.}

\begin{corollary}[Spot checks of a committed training log]\label[corollary]{cor:training:log}
Run \cref{prot:training:spot} with $q$ checks, and let $h=\lceil\log_2(N+1)\rceil$. Let a developer run in time $t^*$, let $\tilde L$ be its string of \cref{thm:training:committed}(ii), the log whose checkpoints it opens most often, and let $\gamma^*=\gamma(2t^*+2t_{\V})$.
\begin{enumerate}
\item A developer that commits to a log from $w_0$ to $\hat w$ with $e$ invalid segments, and opens it, is accepted with probability $(1-e/N)^q$.
\item \emph{Work.} The auditor accepts while $\tilde L$ is not a log from $w_0$ to $\hat w$ with fewer than $\varepsilon N$ invalid segments with probability at most $(1-\varepsilon)^q+\sqrt{(2q+2)(N+1)\gamma^*}$.
\item \emph{Result.} If $\hat w$ is not reachable, the developer is accepted with probability at most $(1-1/N)^q+\sqrt{(2q+2)(N+1)\gamma^*}$, and the honest log with $\hat w$ as its last checkpoint is accepted with probability $(1-1/N)^q$.
\item The auditor replays $q$ segments, a fraction $q/N$ of the run, and evaluates $H$ at most $(2q+2)(h+1)$ times; the developer opens at most $2q+2$ checkpoints, with $h$ hashes each.
\end{enumerate}
\end{corollary}

\begin{proof}
The plan is to apply \cref{thm:training:committed} to the spot-check verifier that reads positions $0$, $N$, $j_i-1$, and $j_i$, so $\kappa=2q+2$ and $M=N+1$, and accepts a string if it starts at $w_0$, ends at $\hat w$, and has valid segments at $j_1,\dots,j_q$. It rejects every string that is not a log from $w_0$ to $\hat w$, and accepts a log $L$ from $w_0$ to $\hat w$ with probability $(1-e(L)/N)^q$ (\cref{prop:training:spotcheck}(i)), which gives (i). An accepting run that reads only checkpoints of $\tilde L$ is one in which the verifier accepts $\tilde L$, which has probability at most $(1-\varepsilon)^q$ when $\tilde L$ is not a log from $w_0$ to $\hat w$ with fewer than $\varepsilon N$ invalid segments; \cref{thm:training:committed}(ii) bounds the other accepting runs, which gives (ii). Part (iii) is (ii) with $\varepsilon=1/N$, since every log from $w_0$ to an unreachable $\hat w$ has an invalid segment, and (i) with $e=1$; part (iv) is \cref{thm:training:committed}(iii).
\end{proof}

The same theorem, applied to the verifier of the PCP theorem (\cref{thm:spotchecks:pcp}), is the argument of \textcite{Kilian1992} in the form in which it is usually taught \parencite{Thaler2022}.\index{Kilian's argument}

\begin{corollary}[Kilian's argument]\label[corollary]{cor:training:kilian}
Suppose that collision-resistant families exist. Then every $L\in\cc{NP}$ has an argument with perfect completeness, in which the prover, given a certificate for $x$, commits to a probabilistically checkable proof of length $M$ polynomial in $n=\abs x$, the verifier sends the coins of $t$ independent runs of a verifier that reads $O(1)$ bits, and the prover answers with $O(t\lambda\log n)$ bits. Every prover that runs in time $t^*$ is accepted on $x\notin L$ with probability at most $2^{-t}+\sqrt{O(t)M\gamma(2t^*+2t_{\V})}$, which is $2^{-t}+\negl(\lambda)$ when $t$, $t^*$, and $n$ are polynomial in $\lambda$.
\end{corollary}

\begin{proof}
By \cref{thm:spotchecks:pcp}, $L$ has a verifier that uses $O(\log n)$ coins, reads a constant number $r$ of bits, accepts some proof with probability $1$ when $x\in L$, and accepts every proof with probability at most $\frac12$ when $x\notin L$. Its proofs have polynomial length, since $O(\log n)$ coins select polynomially many positions, and the known constructions compute them from a certificate in polynomial time (\cref{ch:spotchecks}). A verifier that chooses its positions adaptively can read all of the at most $2^r$ positions it might read, a constant number. For a fixed string the $t$ runs are independent, so $\delta(x)\le2^{-t}$, and $t$ runs read $\kappa=O(t)$ positions. \Cref{thm:training:committed} gives the soundness bound; each opening is one bit and $\lceil\log_2M\rceil=O(\log n)$ hashes of $\lambda$ bits.
\end{proof}

What the probabilistically checkable proof adds is distance: a false claim makes every proof fail a constant fraction of the checks, while a false claim about the result of training never needs more than one invalid segment, which $q$ checks find with probability about $q/N$ (\cref{cor:training:log}(iii); \cref{ex:training:backward,ex:training:distance}).

\Textcite{Micali2000} removed the verifier's message: the prover computes the coins itself, as the hash of the root, by the transform of \textcite{FiatShamir1986}.\index{Fiat--Shamir transform}

\begin{theorem}[Non-interactive committed spot checks; \textcite{Micali2000}]\label{thm:training:micali}
In the random-oracle model (\cref{def:toolkit:rom}), let the prover of \cref{thm:training:committed} build its tree with the oracle, derive $\rho$ from the oracle's answers on the root, and send the root and the openings in one message. A prover that makes at most $Q$ queries to the oracle is accepted on $x$ with probability at most $Q\,\delta(x)+O(Q^22^{-\lambda})$.
\end{theorem}

We do not prove the theorem; the bound in this form follows from the analysis of \textcite{BenSassonChiesaSpooner2016}, as presented by \textcite{Thaler2022}. The term $Q\delta(x)$ is harmless for a probabilistically checkable proof repeated until $\delta(x)\le2^{-\lambda}$, and not for a training log: a developer with $\varepsilon N$ invalid segments can change one padding leaf, recompute the root with $h+1$ evaluations, and try again, and about $(1-\varepsilon)^{-q}$ tries suffice\index{grinding} (\cref{ex:training:grinding}). A non-interactive proof of training must check far more segments, or draw its checks from randomness the developer cannot query, such as the times of the hardware snapshots of \cref{sec:training:valid}.

\section{The theory: what a valid log proves}\label{sec:training:valid}

\Cref{cor:training:log} certifies that the developer stands behind a log with few invalid segments. That is only as strong as the specification, and the tolerance weakens it further.

\begin{proposition}[Valid logs without training]\label[proposition]{prop:training:spoof}\index{spoofing}
Consider logs as in \cref{def:training:spec}, with segments of stochastic gradient descent.
\begin{enumerate}
\item \emph{A learning rate that the log reports.} Let $F_j^\eta$ run $\Delta$ steps of stochastic gradient descent with a learning rate $\eta>0$ that the log reports, and let every stochastic gradient have norm at most $G$. For all $\tilde w_0,\hat w\in\R^d$ and $N\ge1$, the interpolation log $W_j=\tilde w_0+\frac jN(\hat w-\tilde w_0)$ satisfies $\norm{F_j^\eta(W_{j-1})-W_j}\le\norm{\hat w-\tilde w_0}/N+\Delta\eta G$ for every $j$. So for every $\tau>0$ the log has no invalid segment once $N\ge2\norm{\hat w-\tilde w_0}/\tau$ and $\eta\le\tau/(2\Delta G)$, and computing it takes $O(Nd)$ arithmetic operations and no gradient.
\item \emph{A fixed learning rate.} Let the specification fix the segment maps, each a contraction with factor $\theta<1$: $\norm{F_j(u)-F_j(u')}\le\theta\norm{u-u'}$ for all $u,u'$. If a log from $w_0$ has no invalid segment, then $\norm{W_j-W^\circ_j}\le\tau(1-\theta^j)/(1-\theta)$ for every $j$, with $(W^\circ_j)$ the honest log, and the bound is attained. For gradient descent with learning rate $\eta$ on quadratic losses whose Hessians have eigenvalues in $[\mu,1/\eta]$, with $\mu>0$, the factor $\theta=(1-\eta\mu)^\Delta$ applies.
\end{enumerate}
\end{proposition}

\begin{proof}
The plan is the triangle inequality for (i) and an induction on the distance to the honest log for (ii). (i) Each of the $\Delta$ steps moves the weights by $\eta$ times a stochastic gradient, so $\norm{F_j^\eta(W_{j-1})-W_{j-1}}\le\Delta\eta G$, while $\norm{W_{j-1}-W_j}=\norm{\hat w-\tilde w_0}/N$; under the two conditions each term is at most $\tau/2$.

(ii) Validity gives $W_j=F_j(W_{j-1})+\xi_j$ with $\norm{\xi_j}\le\tau$, and $W^\circ_j=F_j(W^\circ_{j-1})$. Subtracting, and writing $D_j=\norm{W_j-W^\circ_j}$, the contraction gives $D_j\le\theta D_{j-1}+\tau$, with $D_0=0$ because both logs start at $w_0$, and induction gives $D_j\le\tau(1+\theta+\dots+\theta^{j-1})$. In dimension one, $F_j(u)=\theta u$, $w_0=0$, and $W_j=\theta W_{j-1}+\tau$ attain it. A step of gradient descent on a quadratic loss with Hessian $A_t$ maps a difference $u-u'$ to $(I-\eta A_t)(u-u')$, and the eigenvalues of the symmetric matrix $I-\eta A_t$ lie in $[0,1-\eta\mu]$, so a product of $\Delta$ such matrices has norm at most $(1-\eta\mu)^\Delta$.
\end{proof}

Part (i) is the \emph{infinitesimal-update} attack of \textcite{FangEtAl2023}, who interpolate between a fresh initialization and a stolen model, report a learning rate close to zero, spend one floating-point operation per weight and step, and pass even a verifier that replays every update. \Textcite{ZhangEtAl2022} had also divided the distance between the models equally, but optimized data on which the gradients nearly vanish, so that each replay ends within the tolerance of the next checkpoint (\cref{ex:training:data}); \textcite{FangEtAl2023} found that this attack fails when the verifier changes the checkpoint interval and tightens the tolerance. They also hid invalid updates below the largest ones, which are the ones checked (\cref{ex:training:topq}), as \textcite{JiaEtAl2021} had foreseen, proposing random checks among the remedies. Part (ii) is the soundness that (i) lacks: when the learning rate, the data, and the length are fixed and segments contract, a log valid within $\tau$ ends within $\tau/(1-\theta)$ of the honest result. Networks are not contractions: replayed in one piece from the same initialization and batches, the residual networks of \textcite[Table~I]{JiaEtAl2021} ended $0.82$ to $0.96$ times as far from the original as two networks trained with different batches are from each other, while step-by-step replays erred by at most $0.001$ to $0.002$ of that distance.

Four assumptions restore part of the soundness: \emph{exact replay}, $\tau=0$, by training at higher precision and logging rounding decisions \parencite{SrivastavaAroraBoneh2024} or by fixing the order of floating-point operations \parencite{ArunEtAl2025}; \emph{committed randomness and data}, as when the seeds of the initialization and of the data order are derived from hashes of the data \parencite{ChoiShavitDuvenaud2023}, with data commitments and timestamps \parencite{FangEtAl2023}; \emph{fixed hyperparameters and length}, which remove the freedom of part (i); and \emph{hardware snapshots} at times the developer cannot predict \parencite{Shavit2023}. Even with all four, \cref{cor:training:log} counts invalid segments, and one invalid segment can carry the run anywhere: the exact result needs a proof of the whole computation, as in \cref{ch:delegation}, or a second trainer.

\begin{protocol}{Refereed training}\label{prot:training:referee}\index{refereed delegation}
\item[\msg{P_1}{V}] The hash $y_1$ of its final checkpoint; each trainer has trained by the specification and hashed every checkpoint.
\item[\msg{P_2}{V}] The hash $y_2$ of its final checkpoint.
\item[\act{V}] If $y_1=y_2$, output $y_1$ and stop. Otherwise set $(a,b)=(0,T)$ and $g_a=H(w_0)$.
\phase{Round, while $b-a\ge2$}
\item[\msg{V}{P}] The midpoint $m=\lfloor(a+b)/2\rfloor$, to both trainers.
\item[\msg{P_1}{V}] The hash $u_m$ of its checkpoint $m$.
\item[\msg{P_2}{V}] The hash $u'_m$ of its checkpoint $m$.
\item[\act{V}] If $u_m=u'_m$, set $a=m$ and $g_a=u_m$; otherwise set $b=m$.
\phase{The final step}
\item[\msg{P_1}{V}] Its checkpoint $a$ and the examples of batch $b$, with their paths to the data's commitment.
\item[\msg{P_2}{V}] The same.
\item[\act{V}] Keep the openings that pass their checks ($H(W)=g_a$ for a checkpoint); if two differ, output the collision and stop. Replay step $b$ from the checkpoint on the batch, and output $y_1$ or $y_2$ for the trainer whose hash at $b$ equals the hash of the result.
\end{protocol}

Here the segments are single steps, $N=T$. Each trainer, $\Prov_1$ or $\Prov_2$, may be the developer or an auditor that repeats the run; the referee $\V$ knows the specification, which includes a key for the hash family, drawn at random, and the root of a Merkle tree over the data; a trainer that fails to answer loses.\sidenote{A trainer that keeps every $\sqrt T$th checkpoint replays at most $\sqrt T$ steps to open one.}

\begin{theorem}[Refereed training]\label{thm:training:refereed}
Let the steps $F_1,\dots,F_T$ be deterministic, let $\tau=0$, and let one of the two trainers be honest: it trains by the specification and answers with the hashes of its checkpoints $W_t=F_t(W_{t-1})$, $W_0=w_0$, and with true openings.
\begin{enumerate}
\item \emph{Correctness.} For every strategy of the other trainer that runs in time $t^*$, the referee outputs $H(W_T)$ except with probability at most $\gamma(t^*+t_{\mathrm{hon}}+t_{\V})$ over the key, where $t_{\mathrm{hon}}$ and $t_{\V}$ are the running times of the honest trainer and of the referee.
\item \emph{Costs.} There are at most $\lceil\log_2T\rceil$ rounds. The referee receives two hashes at the start and two per round, two checkpoints, and the examples of one batch with their paths; it replays one step and hashes at most four checkpoints, besides checking the paths of the batch. The honest trainer runs $T$ steps and hashes $T+1$ checkpoints.
\end{enumerate}
\end{theorem}

\begin{proof}
The plan is to follow two facts through the bisection: the hash $g_a$ is the honest trainer's hash of its checkpoint $a$, and the two trainers' hashes of checkpoint $b$ differ. Before the first round they hold, because the referee computes $g_0=H(w_0)$ and because $y_1\ne y_2$; otherwise the referee outputs the common value, which is the honest $H(W_T)$. In a round where the hashes of checkpoint $m$ agree, the other trainer's hash equals the honest one, so the new $g_a$ is the honest hash. In a round where they differ, the new $b$ is a point of disagreement. The length $b-a$ shrinks to at most $\lceil(b-a)/2\rceil$ in each round, so after at most $\lceil\log_2T\rceil$ rounds $b=a+1$, as in the proof of \cref{thm:debate:bisection}(iii).

At the final step the honest trainer opens $W_a$, whose hash is $g_a$, and the true examples of batch $b$, so the referee holds accepted openings. If the other trainer's accepted openings differ from these, then the referee holds two different checkpoints with the hash $g_a$, or two different openings of one example against the data's root, and \cref{thm:training:merkle}(i) turns either into a collision. Otherwise the referee replays $F_b(W_a)=W_b$, by determinism, whose hash is the honest trainer's hash at $b$ and differs from the other's; so it outputs the honest trainer's final hash $H(W_T)$. The referee fails only when a collision appears, and an algorithm that runs both trainers and the referee and outputs that collision runs in time $t^*+t_{\mathrm{hon}}+t_{\V}$, which gives (i). The costs in (ii) can be read off the protocol.
\end{proof}

The theorem is a relative of the bisection debate (\cref{thm:debate:bisection}). There one debater proposes the midpoint and the other chooses the half that contains the error; here both trainers propose a midpoint, the referee chooses the half by comparing them, and the honest trainer plays both honest debaters at once; hashes replace configurations, too large to send in every round. \Textcite{CanettiRivaRothblum2011,CanettiRivaRothblum2013} proved the general form, for Turing machines, with logarithmically many rounds and every collision-resistant hash family.\sidenote{They also adapted it to x86 programs and built a prototype, Quin, for Windows executables \parencite{CanettiRivaRothblum2011}.} \Textcite{SrivastavaAroraBoneh2024} adapt it to training, following the verification game of \textcite{TeutschReitwiessner2019}: an auditor that repeats the run descends the Merkle tree of the trainer's checkpoints until it finds two consecutive checkpoints whose hashes diverge, for a judge that can perform only minimal computations.\sidenote{One of the auditors must be honest, which \textcite{SrivastavaAroraBoneh2024} call a 1-of-$n$ honesty assumption.}

\paragraph{Compute monitoring.}\index{compute monitoring} \Textcite{Shavit2023} proposed that the hardware make the commitments:\sidenote{His Appendix~B, which derives the bound of \cref{prop:training:monitor}, writes the probability of catching the run as $(1-p_n)^{\lceil T/T_m\rceil}$, with $p_n$ the probability of missing it in one period, where $1-p_n^{\lceil T/T_m\rceil}$ is meant; the next line is correct. His Table~1 lists $132$ inspections a year for a run of Chinchilla's size among $10^3$ chips; with $T_m=T=30$ days the bound gives $118$, and $132$ comes from the approximation $\ln(1-x)\approx-x$ (\cref{ex:training:shavit}).} firmware interrupts each chip at random times, with a small probability at every clock cycle, copies the weights it holds to the developer's storage, and keeps their hash where the developer cannot alter it; inspectors sample chips, and a developer whose chips hold snapshots must produce a transcript that explains them. The snapshots are commitments at times the developer cannot predict, so a log cannot be composed after the run.

\begin{proposition}[Detection by inspecting chips]\label[proposition]{prop:training:monitor}
Let a run occupy $C_{\mathrm{run}}$ of a developer's $C$ chips during $p$ monitoring periods of length $T_{\mathrm m}$, and let each chip take snapshots at the times of a Poisson process of rate $f$, independently of the other chips. If in each period the verifier inspects $m\le C$ distinct chips, chosen uniformly and independently of other periods, then some inspected chip holds a snapshot of the run taken in its period with probability at least
\[1-\Bigl(1-\frac{C_{\mathrm{run}}}C\bigl(1-e^{-fT_{\mathrm m}}\bigr)\Bigr)^{mp}.\]
\end{proposition}

\begin{proof}
The plan is to bound the probability of missing the run in one period by a product over the inspected chips, and to multiply over periods. Fix a period, let $\phi=1-e^{-fT_{\mathrm m}}$ be the probability that a chip takes at least one snapshot in it, and let $K_i$ be the number of run chips among the first $i$ inspected. Distinct run chips hold snapshots independently, so the period is missed with probability $\E[(1-\phi)^{K_m}]$. We show by induction that $\E[(1-\phi)^{K_i}]\le(1-\phi C_{\mathrm{run}}/C)^i$; it holds for $i=0$. Given $K_{i-1}$, the $i$th chip is a run chip with probability $(C_{\mathrm{run}}-K_{i-1})/(C-i+1)$, so, with $g=(1-\phi)^{K_{i-1}}$ and $k=C_{\mathrm{run}}-K_{i-1}$, we have $\E[(1-\phi)^{K_i}]=\E[g]-\phi\,\E[gk]/(C-i+1)$. Both $g$ and $k$ are nonincreasing functions of $K_{i-1}$, so $\E[gk]\ge\E[g]\E[k]$: for two independent copies of $K_{i-1}$ the product of the differences of $g$ and of $k$ is nonnegative, and its mean is $2\E[gk]-2\E[g]\E[k]$. Each inspected chip is a run chip with probability $C_{\mathrm{run}}/C$, so $\E[k]=C_{\mathrm{run}}(C-i+1)/C$, and therefore $\E[(1-\phi)^{K_i}]\le\E[g](1-\phi C_{\mathrm{run}}/C)$, which completes the induction. Inspections and snapshots in different periods are independent, the latter because a Poisson process has independent increments, so the probabilities of missing the run in the $p$ periods multiply.
\end{proof}

With $C_{\mathrm{run}}/C$ for his $H/(aTC)$ and $p$ for $T/T_m$, this is the bound behind Equation~(1) of \textcite{Shavit2023}, which he solves for the number of inspections; his inspections are independent draws, and distinct ones only help. It assumes that the run uses the fewest chips it can, that every chip carries the firmware and is registered, and that the stored hashes cannot be altered; a detected run must still be explained by a transcript that a trusted cluster replays in segments, which brings back \cref{cor:training:log,prop:training:spoof}.

\section{Running the protocol}\label{sec:training:running}

The script \texttt{14-training-pol.py} trains logistic regression with $d=32$ weights on $4{,}096$ synthetic examples by stochastic gradient descent in single precision (batches of $64$, learning rate $0.5$, $T=1{,}600$ steps, a checkpoint every $\Delta=8$ steps, $N=200$), and runs \cref{prot:training:spot} with keyed SHA-256. To emulate nondeterministic hardware, the auditor's replay adds the per-example gradients of each step in a fresh random order; a segment then replays within $5.5\times10^{-7}$ of its checkpoint, against segment displacements of $0.049$ to $0.25$. Forged logs skip $e=4$ or $e=20$ segments, writing an extrapolation of the previous update instead, which replays at least $0.067$ from its checkpoint; here $\tau=10^{-4}$. The spoofs of \cref{prop:training:spoof}(i) interpolate from a fresh draw of the initialization\sidenote{Their audits take this draw as the specified initialization; proof-of-learning checks only the distribution of the initialization.} to the honest final model, at distance $D=13.2$: spoof~A with $N'=50{,}000$ segments and a reported learning rate of $5\times10^{-6}$, spoof~B with $N'=200$ and the same rate, and spoof~C with $N'=200$ and the specified rate. The script \texttt{14-training-monitor.py} simulates \cref{prop:training:monitor} for a run on $222$ of $1{,}000$ chips, the size of Chinchilla's run on H100 chips in $30$ days in Shavit's accounting, with $f=0.1$ snapshots a day and $p=3$ periods of $10$ days.

\begin{figure}[t]
\begin{fullwidth}
\centering
\begin{tikzpicture}
\begin{groupplot}[group style={group size=3 by 1,horizontal sep=11mm},vplot,
    width=0.345\linewidth,height=0.25\linewidth,legend style={font=\scriptsize}]
\nextgroupplot[xlabel={checks $q$},ylabel={rate of detection},xmin=0,xmax=95,ymin=0,ymax=1.06,
    title={(a) forged logs, $N=200$},legend style={at={(0.97,0.04)},anchor=south east}]
\addplot[theory,forget plot] table[x=q,y=d20]{figures/data/14-training-pol-detect-theory.dat};
\addplot[theory2,forget plot] table[x=q,y=d4]{figures/data/14-training-pol-detect-theory.dat};
\addplot[cheat,ebars] table[x=q,y=d20,y error plus=d20ep,y error minus=d20em]{figures/data/14-training-pol-detect.dat};
\addplot[cheat2,ebars] table[x=q,y=d4,y error plus=d4ep,y error minus=d4em]{figures/data/14-training-pol-detect.dat};
\legend{$e=20$,$e=4$}
\nextgroupplot[xmode=log,xlabel={tolerance $\tau$},ylabel={rate of acceptance},xmin=1e-9,xmax=10,ymin=-0.04,ymax=1.06,
    xtick={1e-8,1e-6,1e-4,1e-2,1},
    title={(b) replay and spoofs, $q=10$},legend style={at={(0.03,0.5)},anchor=west}]
\addplot[theory,forget plot] table[x=tau,y=A]{figures/data/14-training-pol-tol-theory.dat};
\addplot[theory2,forget plot] table[x=tau,y=C]{figures/data/14-training-pol-tol-theory.dat};
\addplot[honest,ebars] table[x=tau,y=h,y error plus=hep,y error minus=hem]{figures/data/14-training-pol-tol.dat};
\addplot[cheat,ebars] table[x=tau,y=A,y error plus=Aep,y error minus=Aem]{figures/data/14-training-pol-tol.dat};
\addplot[cheat2,ebars] table[x=tau,y=C,y error plus=Cep,y error minus=Cem]{figures/data/14-training-pol-tol.dat};
\legend{honest log,spoof A,spoof C}
\nextgroupplot[xlabel={chips inspected per period $m$},ylabel={rate of detection},xmin=0.5,xmax=12.5,ymin=0,ymax=1.06,
    title={(c) compute monitoring},legend style={at={(0.97,0.04)},anchor=south east}]
\addplot[bound,forget plot] table[x=m,y=bound]{figures/data/14-training-monitor-theory.dat};
\addplot[theory,forget plot] table[x=m,y=exact]{figures/data/14-training-monitor-theory.dat};
\addplot[cheat,ebars] table[x=m,y=det,y error plus=detep,y error minus=detem]{figures/data/14-training-monitor.dat};
\legend{simulated runs}
\end{groupplot}
\end{tikzpicture}
\end{fullwidth}
\caption{The protocols run against the theory. (a) The fraction of $1{,}000$ audits per point that reject a log with $e$ of its $200$ segments forged, against the number $q$ of checks; lines: $1-(1-e/N)^q$ (\cref{cor:training:log}(i)). (b) The fraction of $400$ audits with $q=10$ that accept, against the tolerance: the honest log replayed in a random order of summation, and spoofs A and C of \cref{prop:training:spoof}, which run no training; lines: the exact acceptance $(n_\tau/N')^{10}$ of each spoof, with $n_\tau$ the number of its segments that replay within $\tau$. (c) The fraction of $20{,}000$ simulated runs found by inspecting $m$ chips per period; solid: the exact probability; dashed: the bound of \cref{prop:training:monitor}, within $0.0003$ of it. Scripts \texttt{14-training-pol.py} (seeds $1401$--$1411$) and \texttt{14-training-monitor.py} (seed $1440$). All intervals are Clopper--Pearson $95\%$ intervals.}\label{fig:training:runs}
\end{figure}

In \cref{fig:training:runs}(a) the forged logs were caught at the rate $1-(1-e/N)^q$: $56.5\%$ of the audits with $q=40$ caught the log with $4$ forged segments (exact $55.4\%$), and $98.5\%$ the log with $20$ (exact $98.5\%$). Of the $22$ intervals, $19$ contain the exact value, and the other three, all for $e=4$ and $q\le8$, miss it by at most $0.0012$. The honest log passed all $1{,}000$ audits with $q=90$. At $\varepsilon=0.02$, detection with probability $0.95$ takes $q=149$ checks, three quarters of the number of segments.

\Cref{fig:training:runs}(b) shows what the tolerance costs. The honest log, replayed in a random order, was accepted in no audit at $\tau\le5.6\times10^{-8}$ and in every audit from $\tau=5.6\times10^{-7}$. Spoof~A, which computed no gradient, is accepted with probability $1$ once $\tau\ge2.6\times10^{-4}$, within the bound $D/N'+\Delta\eta G=3.2\times10^{-4}$ of \cref{prop:training:spoof}(i), with $G=1.51$ the largest norm of an example. Spoof~B, with $250$ times fewer segments, needs $\tau\ge0.066=D/N'$; spoof~C, which keeps the specified learning rate, passes only when the tolerance reaches the size of a training segment, for its segments replay between $0.043$ and $0.21$ from their checkpoints. All $123$ intervals of the three spoofs contain their exact values. The tolerance that admits the honest replay, about $10^{-6}$, would admit spoof~A with about $10^{7}$ segments, cheap to compute and claiming about $10^{8}$ steps.

\begin{table}[t]
\centering\small
\begin{tabular}{@{}l r r r r r r r@{}}
\toprule
& & \multicolumn{4}{c}{auditor or referee} & \multicolumn{2}{c}{honest trainer}\\
\cmidrule(lr){3-6}\cmidrule(l){7-8}
protocol & $T$ & steps & hashes & kB & ms & steps & ms\\
\midrule
spot checks, $q=10$ & \num{1600} & $80$ & $189$ & $8.1$ & $0.63$ & \num{1600} & $7.1$\\
refereed, $8$ rounds & $2^{8}$ & $1$ & $836$ & $33.9$ & $0.33$ & $256$ & $1.4$\\
refereed, $12$ rounds & $2^{12}$ & $1$ & $836$ & $34.1$ & $0.36$ & \num{4096} & $22$\\
refereed, $16$ rounds & $2^{16}$ & $1$ & $836$ & $34.4$ & $0.42$ & \num{65536} & $348$\\
\bottomrule
\end{tabular}
\caption{Costs of one audit by \cref{prot:training:spot} and of one refereed run of \cref{prot:training:referee}, on the training of \cref{fig:training:runs}. Hashes are SHA-256 evaluations; kB counts the openings and hashes received, for the referee mostly the paths of a batch's $64$ examples. Times, in Python on a laptop, vary between runs. Scripts \texttt{14-training-pol.py} and \texttt{14-training-referee.py}.}\label{tab:training:costs}
\end{table}

The script \texttt{14-training-referee.py} runs \cref{prot:training:referee} with an honest trainer against four cheats: one skips a step, one trains one step on a batch with flipped labels, one stops halfway, and one reports the honest model plus a Gaussian perturbation of scale $0.05$ per weight. Over $T=2^8,\dots,2^{16}$ the cheaters won none of $266$ refereed runs, and every run took exactly $\log_2T$ rounds and one replayed step (\cref{tab:training:costs}). From $T=2^8$ to $2^{16}$ the honest trainer's work grows $256$-fold and the referee's not at all, apart from two hashes received in each extra round; a spot-checking auditor that wanted to catch a single skipped step with probability $\frac12$ would have to replay half the run.

\section{The three questions}\label{sec:training:questions}

\questionone With one developer, the game pays the developer when the auditor accepts. If the specification fixes initialization, data, order, hyperparameters, and length, and replay is within the tolerance, then the honest developer is accepted with probability $1$ (\cref{cor:training:log}(i)), so honest training is a best response, at the cost of the training, $2^{h+1}-1$ hashes, and stored checkpoints. It is enforced only up to the detection probability: skipping $\varepsilon N$ segments saves that much computation at a risk of about $1-(1-\varepsilon)^q$. When the log reports its own learning rate, a spoof that does no training is accepted with probability $1$ (\cref{prop:training:spoof}(i)), and honest training is not an equilibrium of a game that charges for computation. In refereed training the honest trainer wins against every opponent unless a collision is found (\cref{thm:training:refereed}); if the loser forfeits a deposit larger than the cost of training, honest training is a best response to an honest opponent, though not dominant: two trainers who deviate alike agree, and nothing is checked.

\questiontwo The auditor of \cref{prot:training:spot} replays $q$ segments, a fraction $q/N$ of the training ($5\%$ in \cref{tab:training:costs}; about $10^{23}$ operations in the oversight box), evaluates $H$ at most $(2q+2)(\lceil\log_2(N+1)\rceil+1)$ times, and exchanges four messages. The referee of \cref{prot:training:referee} spends $\lceil\log_2T\rceil$ rounds of two hashes and replays one step. A zero-knowledge proof of training for VGG-11, about ten million weights, is checked in about $130$ milliseconds from $1.63$ megabytes, whatever the number of steps, but costs the prover about $15$ minutes per step on a batch of $16$ \parencite{AbbaszadehEtAl2024}. Compute monitoring inspects $m$ chips per period; $6$ give $90\%$ detection in \cref{fig:training:runs}(c).

\questionthree Against a developer that cannot find collisions, a committed log with a fraction $\varepsilon$ of invalid segments is caught with probability at least $1-(1-\varepsilon)^q$, less a negligible term (\cref{cor:training:log}(ii)); $t$ audits of one commitment with fresh checks act as $tq$ checks. A wrong model needs only one invalid segment and is caught with probability about $q/N$ (\cref{cor:training:log}(iii)), and against a log valid within the tolerance there is no guarantee (\cref{fig:training:runs}(b)). Refereed training outputs the honest result unless a collision is found. Compute monitoring finds a run with the probability of \cref{prop:training:monitor}, if every chip carries the firmware.

\Needspace{14\baselineskip}
\protocolcard{a committed log to the model has few invalid segments; refereed: the model is the exact result of the run}{a developer $\Prov$ and an auditor $\V$; refereed: two trainers, $\Prov_1$ and $\Prov_2$, and a referee $\V$}{passes within $\tau$, as does a spoof if $\eta$ is free; refereed: best response to an honest opponent}{$q$ of $N$ segments replayed and $O(q\log N)$ hashes; refereed: $\lceil\log_2T\rceil$ rounds and one step}{against bounded developers, $1-(1-\varepsilon)^q-\negl(\lambda)$ for a fraction $\varepsilon$ of invalid segments, about $q/N$ for a wrong model, none within $\tau$; refereed: $1-\negl(\lambda)$}{collision resistance; data, seed, hyperparameters, and length fixed; refereed: determinism, an honest trainer}

\section{What the model leaves out}\label{sec:training:leaves}

\paragraph{What was idealized away.} The protocol certifies a relation between a log and a specification, and the overseer must obtain a specification that says what matters: a commitment binds the developer to some data, not to where the data came from, whether they are licensed, or whether they carry a planted trigger. Replay is idealized as a function; on real hardware it is a distribution, and the tolerance it forces is where spoofs live. Our logistic regression replays a whole run within $7.8\times10^{-7}$ of the original, the networks of \textcite{JiaEtAl2021} not at all (\cref{sec:training:valid}). Exact replay has a price: in a prototype, trainer overheads of $1.2$ to $1.4$ times, auditor overheads of $1.3$ to $1.7$ times, and about $20$ megabytes of logged rounding decisions per step of GPT-2 \parencite{SrivastavaAroraBoneh2024}. The deployed model must still be tied to the published one (\cref{ch:zeroknowledge}). Replaying a fraction of a frontier run is itself a large run, which \textcite{Shavit2023} places on a trusted cluster, air-gapped so that the developer's data do not leak; and the auditor sees checkpoints and data, which \textcite{ChoiShavitDuvenaud2023} handle by trusting the verifier to keep them secure.

\paragraph{Where it breaks.} The security argument of proof-of-learning is a heuristic, and its authors say so: the cost of forging rests on a conjecture about inverting training \parencite{JiaEtAl2021}. \Textcite{ZhangEtAl2022} spoofed it with an adversary that holds the final weights and may modify the training data, and argued, on CIFAR-10, CIFAR-100, and a subset of ImageNet, that the spoofs cost significantly less than training. \Textcite{FangEtAl2023} gave attacks that work whatever the verifier's interval and tolerance (\cref{sec:training:valid}), found no way to produce a fully valid log more cheaply than by training, and concluded that formally proving the robustness of a verifier for proof-of-learning is not currently possible. The defenses of \textcite{ChoiShavitDuvenaud2023}, a committed seed and tests of memorization, detect all known attacks from the literature on proof-of-learning in their experiments with GPT-2-sized models, but are offered as heuristics, and by their authors' account cannot detect small additions of data, such as an inserted backdoor, or small changes to the weights. A trainer who knows the auditor's hardware can choose among the $2^n$ models that its $n$ borderline rounding decisions allow; \textcite{SrivastavaAroraBoneh2024} argue that this cannot embed an arbitrary property, such as a given backdoor, and leave its full analysis open. Refereed training fails when no trainer is honest or trainers collude, and non-interactive proofs of a training log are open to grinding (\cref{thm:training:micali}). Proofs of the whole computation remain expensive: for logistic regression on $262{,}144$ records with $1{,}024$ features, the prover of \textcite{GargEtAl2023} takes under $10$ minutes in its online phase and close to an hour in a phase that depends on the data.

\paragraph{Open problems.} \Textcite{Shavit2023} defines protocols for training transcripts that are efficient for the verifier and hard to spoof, and leaves provably secure ones open. \Textcite{FangEtAl2023} ask whether spoofing algorithms cheaper than a constant fraction of training exist for networks at all, and which representative subsets of updates a verifier could check; for a class of networks, a positive answer would be a tracking bound like \cref{prop:training:spoof}(ii) for segments that are not contractions. Spot checks with distance, which would turn a wrong final model into many invalid checks without a full probabilistically checkable proof of the run, would close the gap between work and result (\cref{ex:training:distance}). Refereed training in which copying a competitor's hashes is not a free strategy, and proofs about data, such as that a model was not trained on a given document, are open as well.

\begin{exercises}
\item \consolidation (The worked run.) In \cref{tab:training:run}, compute the probability that the lazy log survives $q=3$ distinct checks, and the least number of independent checks that catches it with probability $0.99$. For the one-jump log of the text, find the least such $q$, and explain why it is smaller than $\lceil8\ln100\rceil$.\label{ex:training:run}

\item \consolidation (Refereed training by hand.) Let the honest and the lazy developer of \cref{tab:training:run} play \cref{prot:training:referee}, each checkpoint serving as its own hash. Which midpoints are queried, which step is replayed, and why does it suffice that the bisection ends at step $7$, although the lazy log first goes wrong at step $5$?\label{ex:training:referee}

\item \consolidation (The tracking bound.) (a) For gradient descent with $\eta=\frac12$ on the one-dimensional loss $\frac{\mu}2w^2$ with $\mu=0.2$ and $\Delta=8$, compute $\theta$ and the bound of \cref{prop:training:spoof}(ii). (b) Show that without the contraction, for $F_j(u)=2u$, a log without invalid segments can end $\tau(2^N-1)$ from the honest one.\label{ex:training:tracking}

\item \consolidation (Shavit's table.) Solve the bound of \cref{prop:training:monitor} for the number $m$ of inspections per period that gives detection probability $0.9$, and evaluate $365m/30$ for $C=10^3$, $10^5$, and $10^7$ chips, with $C_{\mathrm{run}}=222$, $f=0.1$ per day and $T_{\mathrm m}=T=30$ days, so that $p=1$. Compare with Shavit's $132$, $1.33\times10^4$, and $1.33\times10^6$, and with what the approximation $\ln(1-x)\approx-x$ gives.\label{ex:training:shavit}

\item \variation (More than two trainers.) Let $n$ trainers run one specification, at least one of them honest. Design a referee that outputs the honest hash after at most $n-1$ runs of \cref{prot:training:referee}. Why can it not take the majority hash?\label{ex:training:many}

\item \variation (A replay that is a distribution.) Suppose each replay of an honest segment exceeds $\tau$ with probability $\beta$, independently, while each invalid segment always does. The auditor accepts if at most $z$ of its $q$ checks fail. Compute its completeness and its soundness error against a log with a fraction $\varepsilon$ of invalid segments, and find $q$ and $z$ that make both errors at most $0.05$ when $\beta=0.01$ and $\varepsilon=0.1$.\label{ex:training:noisy}

\item \attack\hard (Grinding.) In \cref{thm:training:micali}, a developer with $\varepsilon N$ invalid segments changes one padding leaf, recomputes the root, and tries again until the checks derived from the root miss the invalid segments. Show that the expected number of tries is $(1-\varepsilon)^{-q}$, evaluate it for $\varepsilon=0.02$ and $q=50$, and find how large $q$ must be for success probability at most $2^{-128}$ against a developer that makes $2^{64}$ queries. Compare with $N=200$.\label{ex:training:grinding}

\item \attack (Checking the largest updates.) Let a log record every step, and let a verifier replay, in each epoch of $S$ steps, the $Q$ steps with the largest updates, as \textcite{JiaEtAl2021} propose. Construct a log in which only $Q$ steps per epoch are valid and which this verifier accepts, as in the second attack of \textcite{FangEtAl2023}. With what probability would $q$ uniform checks of the whole log catch it?\label{ex:training:topq}

\item \attack (Data the developer chooses.) Let each step use one example $(x,y)$ and the loss $\frac12(\ip{x,w}-y)^2$ with the specified learning rate $\eta$. Show that for all $u\ne v$ in $\R^d$ some example satisfies $u-\eta(\ip{x,u}-y)x=v$, so that every interpolation log is valid with $\tau=0$ if the developer chooses the data, the idea behind the attacks of \textcite{ZhangEtAl2022}. What does committing to the data before training prevent?\label{ex:training:data}

\item \attack\hard (Backward chaining.) For gradient descent on a quadratic loss whose Hessian has eigenvalues in $(0,1/\eta)$, with $\tau=0$, show that each step is invertible, and that a developer with an arbitrary model $\hat w$ can build a log from $w_0$ to $\hat w$ whose single invalid segment lies at a position of its choice, at the cost of about $N$ segments. Conclude that $q$ uniform checks then detect it with probability at most $q/N$, and that always checking the last segment does not help.\label{ex:training:backward}

\item \simulation (How long must a spoof be?) In \texttt{14-training-pol.py}, build spoof~A with $N'=10^2,\dots,10^5$ segments and reported learning rates $10^{-4}$ and $10^{-6}$ times the specified one, and plot against $N'$ the least tolerance at which it is always accepted. Compare with $D/N'+\Delta\eta G$: which term dominates, and what follows for a specification that fixes the number of steps?\label{ex:training:spooflength}

\item \research (Spot checks that see the result.) Find an encoding of a training run, computable at a small multiple of the cost of training, in which every log that ends at a model other than the honest result fails a constant fraction of the checks. Probabilistically checkable proofs do this for every computation, at polynomial cost and with proofs of quasilinear length in the best constructions (\cref{ch:spotchecks}); whether gradient descent admits a cheaper encoding is open.\label{ex:training:distance}
\end{exercises}

\notesandrefs
Hash trees are from \textcite{Merkle1987}, in a digital signature built from a conventional encryption function. \Textcite{Kilian1992} used them to commit to probabilistically checkable proofs, which gave arguments for $\cc{NP}$ with polylogarithmic communication; \textcite{Micali2000}, announced in 1994, made such arguments non-interactive with a random oracle, by the transform of \textcite{FiatShamir1986}; and \textcite{BenSassonChiesaSpooner2016} extended the construction to interactive oracle proofs. \Textcite{Thaler2022} gives a textbook account. Refereed delegation is due to \textcite{CanettiRivaRothblum2011,CanettiRivaRothblum2013}, and \textcite{ArunEtAl2025} apply it to machine-learning programs.

\Textcite{JiaEtAl2021} defined proof-of-learning for disputes about the ownership of a model and for detecting faulty updates in distributed training; \textcite{ChoiShavitDuvenaud2023} test training data at a cost they put as low as $1.3\%$ of the training compute, besides some replay; and \textcite{Shavit2023} put these tools into a design for verifying rules on large training runs, from chip firmware to the supply chain. \Textcite{GargEtAl2023} formalized zero-knowledge proofs of training and built one for logistic regression, \textcite{AbbaszadehEtAl2024} reached a network with ten million weights, and \textcite{WaiwitlikhitEtAl2024} prove training by stochastic gradient descent in fixed point inside audits that reveal neither data nor weights (\cref{ch:zeroknowledge}).

\paragraph{Translating the literature.} \Textcite{JiaEtAl2021} write $W_t$ for the weights after step $t$, $k$ for our checkpoint interval $\Delta$, $Q$ for the number of updates replayed in each epoch, chosen by size rather than at random like our $q$, $\delta$ for our tolerance $\tau$, $E$ for the number of epochs, and $S$ for the steps per epoch, so that $T=ES$; \textcite{FangEtAl2023} follow them. \Textcite{Shavit2023} writes $H$ for the computation of a run in floating-point operations, $a$ for a chip's operations per day, $T$ for the length of a run in days, $C$ for the number of chips, $s$ for the chips inspected per period (our $m$), $T_m$ for the period, and $f$ for the rate of snapshots; his $H/(aT)$ is our $C_{\mathrm{run}}$ and his $T/T_m$ our $p$. \Textcite{CanettiRivaRothblum2011} call the referee the client and the trainers servers; \textcite{SrivastavaAroraBoneh2024} call them a judge, a trainer, and an auditor.

\printchapterbib
\end{refsection}
